\section{Public Key Infrastructure}

Public key cryptography is used in many applications:
\begin{itemize}
    \item Secure web connections
    \item Software signing (against malware)
    \item Secure messaging and email
    \item Cryptocurrencies and blockchains
\end{itemize}

How do we know the public key of an entity?
Through digital certificates!

A digital certificate is a signature binding an entity to a public key
$\text{cert} = \text{Sign(“Bob's key is $pk_B$”)}$.
Who will sign this?
Once a single public key is securely transferred,
we use it to trust other's keys.
Specifying such details and other things
(roles, policies, software, etc) constitute a
public key infrastructure.
The simplest PKI consists of a single certificate authority (CA).
A certificate authority is trusted to issue certificates
for everyone's public keys. Of course,
it's ot very practical for everyone to trust the same CA.
One CA'ss certification process may be insufficient.
The CA is a single point of failure, but luckily
we can rely on multiple CAs.
Bob decides which CAs he wants to issue a certificate.
Alice, presented with multiple certificates, decides
which ones she trusts.
The security is as good as the least-secure CA that Alice trusts.
Another approach is through certificate chains.

A “root CA” can delegate issuing certificates to other parties.
Each of these CAs can then issue certs for other parties.

What if someone's private key is leaked who previously got a certificate for their
public key?
We need to be able to generate a new private public pair and remove the old
public key from circulation.
\begin{itemize}
    \item Expiration: have certificates expire after a certain date:
    \item Revocation: have certificates include a serial number, and revoke certificates
          based on that.
\end{itemize}

To visit a website securely, each website needs to obtain a host 
certificate from a CA.
The certificate includes info such as the domain name, 
person or organization it
was issued for, issue date and expiration date of the 
certificate, and the public key.
The user can verify that the website admin owns the website name 
(which prevents spoofing).
The website encrypts the connection between user browser and itself 
(which prevents eavesdropping).

PGP uses hybrid encryption
(symmetric-key + public-key).
A symmetric key is generated by the sender.
The symmetric key is used to encrypt the
message, and is used only once (aka
session key).
The session key is sent over using
receiver's public key.
Web of trust is a decentralized
alternative to hierarchical CAs.
Every user signs certificates for whoever they
want.